%!TEX root = ../dokumentation.tex
\chapter{Projektbeschreibung}
\label{chap:projektbeschreibung}

\section{Hintergrund}

Die Qualifizierung und Weiterentwicklung der Mitarbeitenden stellt für
dm-drogerie markt einen zentralen Erfolgsfaktor dar. Als international
tätiges Handelsunternehmen mit rund 90.000 Mitarbeitenden entstehen
jährlich umfangreiche Aktivitäten im Schulungs- und
Weiterbildungsmanagement, beispielsweise bei der Planung von
Schulungen, der Verwaltung von Teilnehmenden sowie der Dokumentation
von Lernfortschritten.

Die derzeitige Organisation dieser Prozesse basiert auf mehreren
historisch gewachsenen Systemen. Dadurch entstehen fragmentierte
Datenstrukturen, redundante administrative Tätigkeiten sowie
eingeschränkte Transparenz über Weiterbildungsaktivitäten und
Kompetenzentwicklungen.

Im Business Case wurde daher die Einführung einer zentralen
Lernplattform als Lösungsoption identifiziert. Ziel ist es,
Schulungsprozesse, Lernangebote und Kompetenzdaten konzernweit
in einer einheitlichen Plattform zu integrieren und damit eine
konsistente technologische Grundlage für ein effizientes und
datenbasiertes Weiterbildungsmanagement zu schaffen.

\section{Problemstellung / Chance}

Das bestehende Schulungsmanagement ist derzeit mit einem hohen
administrativen Koordinationsaufwand verbunden. Ein großer Teil der
Tätigkeiten entfällt auf manuelle Planungs-, Abstimmungs- und
Dokumentationsprozesse.

Der jährliche administrative Aufwand wird auf rund
10 Millionen Euro geschätzt. Gleichzeitig erschwert die Nutzung
mehrerer Systeme eine konsistente Aggregation von
Weiterbildungs- und Kompetenzdaten.

Die Einführung einer zentralen Lernplattform ermöglicht die
Standardisierung von Prozessen, die Automatisierung administrativer
Tätigkeiten sowie eine zentrale Datenhaltung. Dadurch können sowohl
Effizienzsteigerungen im Schulungsmanagement als auch eine konsolidierte
Datenbasis für Reporting und Kompetenzmanagement erreicht werden.

\section{Überblick über die angestrebte Wirkung}

\begin{table}[H]
\centering
\begin{tabular}{p{4cm} p{10.5cm}}
\toprule
\textbf{Wirkungsbereich} & \textbf{Beschreibung} \\
\midrule
Strategisch &
Stärkung der Kompetenzentwicklung der Mitarbeitenden und Unterstützung
der strategischen Personalentwicklung. \\

\addlinespace
Finanziell &
Reduktion administrativer Aufwände sowie Konsolidierung bestehender
Systemlandschaften. Die Wirtschaftlichkeitsbetrachtung zeigt eine
sehr kurze Amortisationsdauer, einen hohen \ac{ROI} sowie einen
deutlich positiven Kapitalwert. \\

\addlinespace
Operativ &
Effizientere Planung, Durchführung und Dokumentation von
Schulungsmaßnahmen durch integrierte End-to-End-Prozesse. \\

\addlinespace
Mitarbeitende / Kunden &
Verbesserter Zugang zu Lernangeboten und höhere Transparenz über
Qualifizierungs- und Entwicklungsmöglichkeiten. \\

\addlinespace
Sonstiges &
Aufbau einer konsolidierten Datenbasis für organisationsweite
Weiterbildungs- und Kompetenzkennzahlen. \\

\bottomrule
\end{tabular}
\end{table}

\section{Finanzielle Betrachtung der Plattformeinführung}

Die wirtschaftliche Bewertung basiert auf folgenden Annahmen:

\begin{table}[H]
\centering
\begin{tabular}{p{8cm} p{5cm}}
\toprule
\textbf{Parameter} & \textbf{Wert} \\
\midrule
Administrativer Schulungsaufwand p.\,a. & 10.000.000 € \\
Effizienzsteigerung & 50\,\% \\
Projektkosten (inkl. Migration) & 500.000 € \\
Betrachtungszeitraum & 5 Jahre \\
Diskontierungszins & 5\,\% \\
\bottomrule
\end{tabular}
\end{table}

Neben der Prozessautomatisierung entstehen Einsparungen durch die
Konsolidierung der bestehenden Systemlandschaft. Gleichzeitig fallen
laufende Kosten für den Betrieb der neuen Plattform an.

\begin{table}[H]
\centering
\begin{tabular}{p{8cm} p{3cm} p{3cm}}
\toprule
\textbf{Kostencluster} & \textbf{Kosten p.\,a.} & \textbf{Einsparung p.\,a.} \\
\midrule
SuccessFactors Lizenzkosten & 1.350.000 € & -- \\
Zusatzspeicher & 13.000 € & -- \\
DEV + TST Tenants & 20.000 € & -- \\
Content Storage & 7.500 € & -- \\

\midrule

Lizenzen, Projekte und Schulungsaufwand bisheriger Systeme & -- & 510.000 € \\
Rechenzentrum und Infrastruktur & -- & 190.000 € \\
Betrieb und Wartung & -- & 195.000 € \\
Datenmanagement und Schnittstellenbetrieb & -- & 75.000 € \\
Ausfall- und Qualitätskosten & -- & 25.000 € \\

\midrule
\textbf{Summe} & \textbf{1.390.500 €} & \textbf{995.000 €} \\
\bottomrule
\end{tabular}
\end{table}

Die zusätzliche Belastung aus Plattformbetrieb beträgt damit

\begin{align}
K_{\text{technologie}} &= 1.390.500 - 995.000 \\
                       &= 395.500\ \text{€ p.\,a.}
\end{align}

Die Reduktion des administrativen Aufwands beträgt

\begin{align}
S_{\text{administrativ}} &= 10.000.000 \cdot 0{,}50 \\
                         &= 5.000.000\ \text{€ p.\,a.}
\end{align}

Daraus ergibt sich eine jährliche Nettoeinsparung von

\begin{align}
S_{\text{netto}} &= 5.000.000 - 395.500 \\
                 &= 4.604.500\ \text{€ p.\,a.}
\end{align}

\paragraph{Payback}

\begin{align}
\text{Payback} &= \frac{500.000}{4.604.500}
\approx 0{,}11\ \text{Jahre} \quad (\approx 1{,}3\ \text{Monate})
\end{align}

\paragraph{ROI}

Über einen Betrachtungszeitraum von 3 Jahren ergibt sich:
\begin{align}
\ac{ROI} &= \frac{3 \cdot 4.604.500 - 500.000}{500.000} \\
         &\approx 26{,}63 \quad (2\,663\,\%)
\end{align}

\paragraph{NPV}

\begin{align}
\ac{NPV} =
\sum_{t=1}^{5}\frac{4.604.500}{(1{,}05)^t}
- 500.000
\approx
\boxed{19{,}4\ \text{Mio.\ €}}
\end{align}

Die Investition amortisiert sich damit innerhalb weniger Monate
und erzeugt über den Betrachtungszeitraum einen deutlich positiven
Kapitalwert.

Darüber hinaus wird das definierte \ac{MOV}-Ziel einer Reduktion
des administrativen Aufwands um 30\,\% deutlich übertroffen:

\begin{equation}
\frac{4.604.500}{3.000.000} \approx 153\,\%
\end{equation}